% arXiv submission kit: Paper 1 (cs.CR primary, cs.LG secondary)
% Converted from docs/arxiv/PAPER1_DRAFT.md (final-form v0, 2026-06-09).
% \pdfoutput=1 tells arXiv's AutoTeX to produce a PDF directly (required
% within the first 5 lines when submitting pdflatex sources).
\pdfoutput=1
\documentclass[11pt]{article}

\usepackage[margin=1in]{geometry}
\usepackage{amsmath}
\usepackage{amssymb}
\usepackage{amsthm}
\usepackage{booktabs}
\usepackage{iftex}
% Content-neutral packaging fix: under an XeTeX-based engine (e.g. tectonic,
% used for local proof-of-compile), select hyperref's xetex driver so it
% matches the real backend. On pdfTeX (arXiv AutoTeX) \ifxetex is false and
% this line is a no-op; the pdfTeX driver loads exactly as before. Nothing
% scientific or typographic in the rendered paper changes.
\ifxetex\PassOptionsToPackage{xetex}{hyperref}\fi
\usepackage[hidelinks]{hyperref}

\theoremstyle{definition}
\newtheorem{algo}{Algorithm}

\title{Chain-of-Custody for AI-Transformed Text: Signed, Replayable,
Distribution-Free Receipts for What a Transformation Preserved}
\author{Umar Syed\\ Independent Researcher\\ \texttt{ototao@pm.me}}
\date{}

% NOTE (kit author): the Markdown draft's header note ("Draft (final-form v0)
% --- 2026-06-09 ... drafting notes are at the foot.") and its closing
% "Drafting notes for the operator" block are operator-facing meta-commentary,
% not paper content; they are preserved as comments at the end of this file.

\begin{document}

\maketitle

\begin{abstract}
Two questions about AI-transformed text lack a portable, offline-verifiable
answer: \emph{who transformed this, and what did the transformation
preserve?} Provider disclosure (EU AI Act Article 50) and image-centric
content provenance (C2PA, SynthID) do not cover text that has been
paraphrased, summarized, or translated --- a manifest detaches on copy, a
watermark is defeated by rewriting. We present a receipt family that answers
both questions for text. A signed, offline-verifiable receipt attests a
transformation (Ed25519 over RFC 8785 JCS-canonical bytes, detached JWS,
JWKS keys); on top of it, a \textbf{distribution-free, replayable
certificate} bounds the expected \emph{meaning-loss} of the transformation
under a \textbf{named judge}. The certificate replays offline over a
committed integer loss vector --- a third party re-runs the conformal
certifier and reproduces the bound to the bit --- while the proof boundary
stays explicit: it bounds a named proxy \emph{marginally}, over an
\emph{i.i.d.} calibration sample and only where that sample matches
deployment, never per-document truth and never ``meaning'' itself. We demonstrate on two public-domain corpora: certified expected
meaning-loss $\le 0.646$ (95\%) for abstractive summarization of US
Congressional bills (BillSum, CC0; $n=64$) and $\le 0.413$ for EN$\to$FR
translation (opus-100; $n=64$), with 39/64 faithful translations scoring
\emph{exactly zero} meaning-loss (under a binary entailment judge at a 0.5
cut) despite near-zero lexical overlap --- the property no watermark or
lexical scheme can certify. The thesis is \textbf{attest, don't detect}: a
signature survives an adversary with a thesaurus; a statistical
``is-this-AI'' classifier does not.
\end{abstract}

\section{Introduction}\label{sec:intro}

Pressure on AI-generated and AI-transformed text is converging on two needs.
First, \textbf{disclosure}: Article 50(2)/(4) of the EU AI Act (applicable
2 August 2026; a grace window to 2 December 2026 for the machine-readable
marking obligation on pre-existing systems is provisionally agreed under the
Digital Omnibus) requires machine-readable, robust marking of AI-generated
content, and the associated \emph{Code of Practice on Transparency of
AI-Generated Content} (finalized 10 June 2026) asks that marking be
``effective, interoperable, robust, and reliable as far as technically
feasible''~\cite{euaiact}. Second, and far less served,
\textbf{accountability for transformation}: when a document is summarized,
re-leveled for a different audience, or translated, \emph{what survived the
operation, and can it be proven to anyone, offline?}

The dominant provenance tools answer neither well for text.
C2PA~\cite{c2pa} binds a cryptographic manifest to a media asset; its 2.4
text binding is a \emph{byte-exact} hard hash, so the manifest verifies the
exact bytes and breaks the moment text is edited, re-flowed, or paraphrased
--- it says nothing about what a transform preserved.
SynthID-Text~\cite{dathathri2024} and statistical ``AI detectors'' target
\emph{generation}, not transformation-preservation, and degrade under
exactly the rewriting text invites --- every general detector has failed
under paraphrase, distribution shift, or bias against non-native writers,
and a major provider retired its own classifier.

We take a different stance: \textbf{attest, don't detect.} Rather than infer
whether text is AI-generated, we let any participating transformer
\emph{attest} what it did and what it preserved, in a receipt anyone
verifies offline --- robust to rewriting because a signature is indifferent
to surface form. The technical contribution is a \emph{composition}: a
signed transformation receipt carrying a \textbf{distribution-free
certificate of meaning-loss under a named judge}, made \textbf{replayable}
so the certificate is reproducible rather than merely asserted. The slogan
is deliberately narrow: what a receipt takes custody of is a
\emph{transformation} and a \emph{named, bounded proxy} for what it
preserved --- \textbf{custody of a transform, not of meaning}. The gap
between ``a bounded named proxy'' and ``meaning'' is never rhetorically
closed; \S\ref{sec:boundary} holds that line in the wire format. Concurrent
work (\S\ref{sec:related}) independently arrived at the same cryptographic
substrate for LLM-API provenance but explicitly disclaims semantic
correctness; our contribution is precisely that disclaimed delta, turned
into a finite-sample, replayable certificate.

\paragraph{Contributions.}
(1) A unified, cross-runtime-verifiable receipt family for
text-transformation provenance. (2) A \emph{meaning-risk certificate}: a
signed, same-commit-replayable, distribution-free upper bound on expected
meaning-loss under a named judge, with its proof boundary enforced in the
wire format. (3) A group-conditional (``perspective'') extension with a
Bonferroni simultaneous guarantee. (4) Two real, audited, public-domain
demonstrations and an empirical coverage validation. (5) An explicit threat
model and the two load-bearing verification preconditions (trusted
out-of-band keys; issuer-attestation vs.\ bound-attestation) that any
``offline-verifiable'' claim must state.

\section{Threat model}\label{sec:threat}

The receipts defend a specific surface; we state it so the guarantees are
not over-read (Table~\ref{tab:threat}).

\begin{table}[htbp]
\centering
\small
\begin{tabular}{@{}p{0.40\textwidth}p{0.21\textwidth}p{0.30\textwidth}@{}}
\toprule
\textbf{Attacker capability} & \textbf{Outcome} & \textbf{Mechanism} \\
\midrule
Tamper with a signed field in transit (bound, corpus, disclosure, \dots) &
\textbf{Rejected} &
Ed25519 over the full JCS-canonical payload \\
\addlinespace
Strip/empty the disclosure (\texttt{not\_covered}/\texttt{disclosure}) &
\textbf{Rejected} &
Stage-A disclosure invariants \\
\addlinespace
Algorithm downgrade (\texttt{alg:none}), \texttt{crit}/\texttt{b64}
downgrade, non-detached JWS &
\textbf{Rejected} &
protected-header pinning \\
\addlinespace
Self-sign a receipt with a \emph{fabricated} bound but an honest loss-hash &
\textbf{Rejected by Stage B} (passes Stage A) &
re-certification over the committed losses \\
\addlinespace
Forge a receipt and present it with an \textbf{attacker-controlled JWKS} &
\textbf{Out of scope} --- verification trusts the JWKS &
keys must come from a trusted root, out-of-band \\
\addlinespace
Replay an expired/stale receipt &
Rejected when a max-age window is enforced &
\texttt{signed\_at} window check \\
\bottomrule
\end{tabular}
\caption{Attacker capabilities, outcomes, and mechanisms.}
\label{tab:threat}
\end{table}

Two preconditions follow and must accompany any ``tamper-evident,
offline-verifiable'' claim: \textbf{(P1)} verification reduces to
\emph{trusting the JWKS}, obtained from a signed trust root out-of-band
(never from the receipt bundle) --- a forgery verifies against an attacker's
own JWKS, as for any JWS system; \textbf{(P2)} the signature attests the
\emph{issuer}, while the conformal \emph{bound} is attested only by the
Stage-B replay (\S\ref{sec:verify}). Out of scope: the honesty of the
issuer's labels (\texttt{scorer}/\texttt{model}/\texttt{provider} are
producer-asserted), the correctness of the transformed output, and any claim
about human-vs-AI authorship.

\section{The substrate and the receipt family}\label{sec:substrate}

All four schemas (\texttt{render}, \texttt{transform},
\texttt{meaning\_risk}, \texttt{perspective}) share one envelope: a payload
signed with \textbf{Ed25519} (RFC 8032~\cite{rfc8032}) over its
\textbf{JCS-canonical bytes} (RFC 8785~\cite{rfc8785}), carried as a
\textbf{detached JWS} (RFC 7515~\cite{rfc7515}) and verified against a
\textbf{JWKS} (RFC 7517~\cite{rfc7517}). The provenance tier
(\texttt{render}/\texttt{transform}) attests \emph{what happened}: a named
transform mapped an input (hash) to an output (hash) under named parameters,
by a named model/provider. The measured-bound tier
(\texttt{meaning\_risk}/\texttt{perspective}) adds the certificate of
\S\ref{sec:certificate}--\ref{sec:verify}.

\paragraph{Cross-runtime determinism.}
A signature is verifiable in another runtime only if that runtime
canonicalizes the payload to \emph{identical bytes}. JCS fixes key ordering
and string/number serialization; the one subtlety is number formatting,
which RFC 8785 \S3.2.2.3 mandates be the ECMAScript
\texttt{Number::toString} form~\cite{ecma262}. We verify Python, Node, and
browser canonicalizers agree byte-for-byte on the receipt payloads,
including across the full sub-$10^{-4}$ float band where a na\"ive
\texttt{repr}-based encoder diverges. The conformal-tier payloads are
additionally \textbf{float-free} --- every rate quantity is an integer
``micro-unit'' (value $\times\, 10^{6}$, $\le 10^{6} \ll 2^{53}$) --- so the
exact-equality replay of \S\ref{sec:verify} is well-defined across runtimes.

\section{The meaning-risk certificate}\label{sec:certificate}

\paragraph{Meaning-loss proxy.}
Fix a \emph{named, versioned} judge
$e:(\text{premise},\text{hypothesis})\to\{0,1\}$ (a local NLI model or
embedding-entailment model; never assumed to equal ``meaning''). For a
(source $S$, output $T$) pair with units $\mathrm{sent}(\cdot)$,
define
\[
\mathrm{recall}=\tfrac{1}{|\mathrm{sent}(S)|}\!\!\sum_{s\in\mathrm{sent}(S)}\!\! e(T,s),\quad
\mathrm{fid}=\tfrac{1}{|\mathrm{sent}(T)|}\!\!\sum_{t\in\mathrm{sent}(T)}\!\! e(S,t),
\]
\[
\ell(S,T)=1-\big(w_r\,\mathrm{recall}+w_f\,\mathrm{fid}\big)\in[0,1],\quad w_r+w_f=1.
\]
Recall penalizes omission, fidelity penalizes fabrication; $\ell(S,S)=0$.
Both committed receipts use $w_r=0.6$ and $w_f=0.4$, weighting omission above
fabrication; the weights are recorded verbatim in each receipt's signed
\texttt{loss\_definition} field, so a verifier reads them rather than
assuming them. The judge is named and versioned in the receipt, so the
certified bound is \emph{explicitly conditional on it}.

The unit $\mathrm{sent}(\cdot)$ in the shipped v1 scorer is a
punctuation-or-newline-delimited segment, not a linguistic sentence:
hard-wrapped sources split at line breaks, and abbreviations split at their
periods. On the BillSum golden 48\% of the 12,149 units begin with a
lowercase character for this reason. The bound is over whatever units the
named scorer produced; the receipt pins the scorer, so replay reproduces the
same units.

\paragraph{Certificate.}
Let $(S_i,T_i)_{i=1}^n$ be a calibration sample drawn i.i.d.\ from the
deployment distribution, $\ell_i=\ell(S_i,T_i)$, preservation
$p_i=1-\ell_i\in[0,1]$. By
Hoeffding's inequality~\cite{hoeffding1963}, the one-sided $(1-\delta)$
lower confidence bound on $\mathbb{E}[p]$ is
$\bar p-\sqrt{\ln(1/\delta)/(2n)}$, hence the $(1-\delta)$ \textbf{upper}
bound on expected meaning-loss is
\begin{equation}
U_\delta \;=\; \min\!\Big(1,\; \bar\ell + \sqrt{\tfrac{\ln(1/\delta)}{2n}}\Big),
\qquad \Pr\big[\mathbb{E}[\ell]\le U_\delta\big]\ge 1-\delta. \tag{1}
\label{eq:bound}
\end{equation}
Three estimators ship: \textbf{Hoeffding} \eqref{eq:bound} for
$\ell\in[0,1]$; \textbf{Clopper--Pearson} (exact) for the binary
preserved/lost view; and the variance-adaptive \textbf{empirical-Bernstein}
bound (Maurer \& Pontil, 2009~\cite{maurerpontil2009}), whose deviation
scales with the sample variance and is therefore tighter for low-variance
batches at larger $n$ (its additive $O(1/(n{-}1))$ term makes it
\emph{looser} than Hoeffding at small $n$ --- a regime fact, not a tuning
knob). The shipped default selects Clopper--Pearson for binary data and
Hoeffding otherwise; the choice is made a priori, never by comparing
realized bounds.

\paragraph{What the precondition actually is.}
Two assumptions carry \eqref{eq:bound} and we separate them, because one is
stronger than this project previously stated. \textbf{Within} the calibration
set the draws must be \emph{independent}: Hoeffding, Clopper--Pearson (an
exact Binomial interval) and the empirical-Bernstein bound are all
inequalities for independent bounded variables, and exchangeability alone
does not suffice. The counterexample is immediate: a sequence that is
all-zeros or all-ones with probability $1/2$ each is exchangeable and bounded
in $[0,1]$ with $\mathbb{E}[\ell]=1/2$, yet $U_\delta$ falls below $1/2$ on
half of all draws, giving coverage $1/2$ against a nominal $0.95$.
\textbf{Between} calibration and deployment, the distributions must match for
the bound to transfer; that assumption is not testable from the calibration
sample and we do not claim it is. Exchangeability is the right precondition
for conformal \emph{prediction}, which is not the procedure used here.
The \texttt{disclosure} strings inside the two committed receipts say
``exchangeability''; that wording is weaker than what \eqref{eq:bound}
requires. We state the stronger requirement here rather than silently re-sign
the receipts, and \S\ref{sec:boundary} carries the same correction.

\paragraph{Wire and disclosure.}
The receipt commits the SHA-256 of the integer-micro loss vector and carries
\texttt{n}, \texttt{delta\_micro}, \texttt{method},
\texttt{point\_estimate\_micro}, \texttt{risk\_upper\_bound\_micro}, the
judge identity, the calibration \texttt{corpus\_id}, and two
\emph{required, enforced} fields: a non-empty \texttt{not\_covered} list
(the layers the proxy structurally cannot reach --- arrangement, sound,
connotation, implicature) and a non-empty \texttt{disclosure} string. A
receipt that discloses nothing cannot pass as a bare bound.

\section{Verification: two stages}\label{sec:verify}

Verification separates \emph{issuer-attestation} (everywhere) from
\emph{bound-attestation} (deterministic, offline, where the certifier runs).

\begin{algo}[{\texttt{verify(envelope, jwks, [losses], [max\_age])}}]
\label{alg:verify}
% Long underscored identifiers cannot hyphenate; \sloppy keeps them out of the
% margin. algo is a theorem-like environment, so this is scoped to it.
\sloppy
\textbf{Stage A (any runtime).} (a) JOSE-verify the detached JWS over
$\mathrm{JCS}(\text{payload})$ against \texttt{jwks}; (b) gate
\texttt{schema} (fail closed on unknown); (c) check header invariants
(\texttt{alg=EdDSA}, matching \texttt{kid}, \texttt{b64:false},
\texttt{crit:["b64"]}); (d) enforce disclosure
(\texttt{not\_covered}${}\ne\emptyset$, \texttt{disclosure}${}\ne{}$``'');
(e) if \texttt{max\_age} given, check the \texttt{signed\_at} window.
$\to$ \emph{authentic + self-disclosing.}\\
\textbf{Stage B (when the loss vector is supplied side-band).} (f) recompute
the loss-vector hash and require equality; (g) re-run the certifier of
\S\ref{sec:certificate} over the committed integer-micro losses at the
payload's \texttt{method}/\texttt{delta}; (h) require the reproduced
\texttt{risk\_upper\_bound\_micro}, \texttt{point\_estimate\_micro}, and
\texttt{n} to match by \textbf{exact integer equality}. $\to$ \emph{bound
reproduced.}
\end{algo}

Stage B is what turns ``measured'' into ``measured and independently
reproducible,'' and it is what distinguishes a \emph{tampered-in-transit}
receipt (signature fails) from an \emph{overclaimed-at-issue} one (signature
valid, bound does not replay). The certificate replays on any machine;
\emph{re-deriving} the losses from raw text additionally requires the named
model judge, whose forward pass is machine-pinned (cross-hardware
floating-point drift) --- a fact the receipt discloses. So ``verifiable in
every runtime'' means the signature and bound \emph{replay}, not the
\emph{meaning recomputation}.

\section{Perspective receipts (group-conditional)}\label{sec:perspective}

A marginal bound hides the cohort it is worst on. Given cohort labels
$g_i$, the \textbf{perspective receipt} certifies \eqref{eq:bound} per
declared cohort over only that cohort's losses --- the discrete-covariate
case of group-conditional risk control (Gibbs, Cherian \& Cand\`es,
2023~\cite{gibbs2023}). The marginal and the per-cohort family are
\emph{separate} $(1-\delta)$ statements; with \texttt{simultaneous} set,
each cohort is certified at $\delta/G$ (Bonferroni) so that \emph{all} $G$
cohort bounds hold \textbf{jointly} at $\ge 1-\delta$. Each cohort pays its
own finite-sample radius --- a small cohort gets a wide bound, which is
honest, not a defect.

\section{Empirical demonstration}\label{sec:empirical}

We issue two real receipts over public-domain corpora. Each is committed (a
signed golden, the integer loss vector, a deterministic generator), replays
offline via Algorithm~\ref{alg:verify} Stage B, was \textbf{independently
re-derived to the exact micro-unit} and adversarially audited before
release.

\subsection{Compression (BillSum, CC0)}\label{sec:billsum}

First 64 bills of the BillSum test split (US Congressional bills + reference
summaries; CC0-1.0 as US-government works), local MiniLM-cosine entailment
judge (mean-pooled \texttt{all-MiniLM-L6-v2} cosine at a 0.5 cut). The output
side of each pair is the dataset's own human-written reference summary, not a
model's; see \S\ref{sec:scope}. The abstractive bill$\to$summary transform
certifies \textbf{expected
meaning-loss $\le 0.6455$ at 95\%} ($n=64$, mean $0.4925$), controlled
against an operator-chosen \textbf{0.70 target} --- an illustrative bar set
by the issuer, not a regulator/SLA threshold, so ``controlled: true'' means
``met the bar the issuer picked'', not ``passed an external quality gate''.
Aggressive summarization loses about half the named proxy on average; the
receipt \emph{certifies how much}, it does not claim little was lost.

\subsection{Translation (opus-100)}\label{sec:translation}

First 64 length-aligned EN$\to$FR pairs of opus-100, local multilingual NLI
judge (mDeBERTa-v3-xnli). As with BillSum, the French side is the corpus's own
reference translation, not a system output. The translation transform certifies
\textbf{expected meaning-loss $\le 0.4124$ at 95\%} ($n=64$, mean $0.2594$),
controlled against an operator-chosen 0.50 target (again illustrative, not
an external bar). The distribution is the headline: \textbf{39 of 64
faithful translations score exactly zero meaning-loss despite near-zero
lexical overlap} between English and French. This headline must be read
\emph{at the judge's resolution}: at the 0.5 NLI cut the per-pair loss takes
only five distinct values $\{0.0, 0.4, 0.6, 0.8, 1.0\}$ --- the 64 committed
losses are $39\times 0.0$, $8\times 0.4$, $8\times 0.6$, $2\times 0.8$,
$7\times 1.0$ --- so ``exactly zero'' means \emph{bidirectional entailment
fired above 0.5 in both directions}, not fine-grained perfect preservation;
the split is also length-aligned and short, which lowers difficulty. With
that caveat stated: a lexical or watermark scheme cannot credit a faithful
translation at all --- there is almost no surface to match --- whereas the
named judge can. What the receipt takes custody of is the \emph{translation
transform} and a \emph{bounded named proxy} for what it preserved, issued as
a signed certificate --- custody of a transform, not of meaning itself.

\subsection{Coverage validation}\label{sec:coverage}

The bounds are only meaningful if valid. We measure empirical coverage of
the $(1-\delta)$ upper bound --- the fraction of trials in which the
certified bound does not fall below the true loss rate --- on synthetic data
with known ground truth, each trial drawing $n=64$ i.i.d.\ Bernoulli losses
(the binary preserved/lost view) at the stated true rate, $2\times 10^{4}$
trials; Table~\ref{tab:coverage} reports the result.

\begin{table}[htbp]
\centering
\begin{tabular}{@{}lcccc@{}}
\toprule
method & tl=0.1, $\delta$=.05 & tl=0.3, $\delta$=.05 & tl=0.5, $\delta$=.05 & tl=0.5, $\delta$=.10 \\
\midrule
target & $\ge 0.95$ & $\ge 0.95$ & $\ge 0.95$ & $\ge 0.90$ \\
Hoeffding & 1.000 & 0.997 & 0.992 & 0.984 \\
empirical-Bernstein & 1.000 & 1.000 & 1.000 & 1.000 \\
Clopper--Pearson & 0.961 & 0.972 & 0.970 & 0.919 \\
\bottomrule
\end{tabular}
\caption{Empirical coverage of the $(1-\delta)$ upper bound on synthetic
data with known ground truth ($n=64$, $2\times 10^{4}$ trials). Here
\texttt{tl} is the \emph{true loss} rate of the generating distribution, the
quantity the bound must not fall below; \emph{target} is the nominal coverage
$1-\delta$. A method is valid where its row meets its target.}
\label{tab:coverage}
\end{table}

All estimators clear their target in every regime (Hoeffding and eB
conservatively, Clopper--Pearson near-nominal). The Bonferroni simultaneous
path's \emph{joint} (all-cohorts) coverage was separately validated at 0.958
against a 0.95 target (Clopper--Pearson; $G=3$ cohorts of $n=60$ drawn
Bernoulli at true loss rates $0.10/0.20/0.30$, $\delta=0.05$,
$2\times 10^{3}$ trials, seed 17). At the receipts' $n=64$ and observed variance,
Hoeffding is the tighter honest estimator than empirical-Bernstein (BillSum
0.6455 vs 0.650; translation 0.4124 vs 0.518), confirming the a-priori default
was not a favorable cherry-pick.

\subsection{Honest scope}\label{sec:scope}

Each receipt is certified under its \emph{own} named judge (the two judges
differ), so the two means are not a single comparable grading scale. The
bounds are over the \emph{named, filtered} calibration corpora --- not over
all summarization or all translation; the length-alignment filter on the
translation corpus plausibly biases its proxy \emph{low} relative to the
unfiltered split. The certificate is a \textbf{batch} primitive: a single
document uses a per-document measurement, not the bound.

Two limits of these demonstrations deserve stating plainly rather than being
left for a reader to discover. First, \textbf{neither output side was produced
by a model}. BillSum's outputs are the dataset's human-written reference
summaries and opus-100's are its reference translations, so what is
demonstrated is that the certificate machinery works end to end over real
(source, output) pairs, not that it has been exercised on the output of an AI
system. The mechanism is indifferent to who produced the output --- the
issuer signs whatever pair it certifies --- but a paper about AI-transformed
text should not let a reader assume these pairs were AI-transformed.
Second, \textbf{judge validity is a separate question from bound validity}. A
verified receipt is a cryptographic fact about a named proxy; it is not
evidence that the proxy tracks human judgment. On SummEval the proxy
correlates only modestly with human faithfulness ratings at summary level
(Spearman $\rho$ between $0.267$ and $0.291$), the NLI judge replicates at
$\rho \approx 0.29$ on FRANK, and the embedding judge used in the BillSum
demonstration above is corpus-dependent and falls to near zero on abstractive
FRANK-XSum. The shipped verifier prints this caveat on every verified
meaning-risk verdict; we repeat it here so the flagship number is not read
without it.

\section{Proof boundary}\label{sec:boundary}

A verified receipt \textbf{proves} (cryptographic, given P1--P2): the named
issuer signed this exact payload; the committed losses hash to the anchor;
the named certifier reproduces the bound on those losses. It \textbf{does
not prove}: output truth or freshness; issuer honesty; that \emph{meaning}
was preserved (only that a \emph{named proxy} is bounded, marginally, over an
i.i.d.\ calibration sample and only where that sample matches deployment, per
\S\ref{sec:certificate}); or anything about human-vs-AI authorship --- we ship no
detection number, and any such signal is advisory, never a guarantee. The
proxy's structural blind spots are declared in the enforced
\texttt{not\_covered} field. This is the discipline that keeps the gap
between ``a bounded named proxy'' and ``preserved meaning'' continuously
visible rather than rhetorically closed.

\section{Position versus prior and concurrent work}\label{sec:related}

\begin{itemize}
\item \textbf{Signed-but-semantics-disclaiming receipts.}
  \textbf{AEX}~\cite{aex2026} uses JCS + SHA-256 + Ed25519 signed
  transformation-receipt chains for LLM APIs but \textbf{explicitly
  disclaims} that an accepted transform is ``semantically correct,
  reasonable, or minimal.'' The same pattern --- a signed artifact that
  proves \emph{what was computed} and avoids any quality claim --- recurs in
  deterministic-inference attestation (EigenAI~\cite{eigenai2026}). Our
  contribution is precisely that disclaimed semantic delta, as a
  finite-sample, replayable certificate.
\item \textbf{The receipt scheme that does not disclaim semantics.} Tool
  receipts for agents~\cite{toolreceipts2026} are the closest counterexample
  to the bullet above, and we name them as one rather than list them among
  the disclaiming schemes: that work cross-references per-claim receipts
  specifically \emph{to detect hallucination}, and reports detection rates by
  error type. Two differences are load-bearing. Its receipts are
  HMAC-authenticated, a symmetric MAC, so they are unforgeable by the model
  under its stated threat model but carry no third-party verifiability and no
  non-repudiation: anyone who can verify can also mint. And its semantic claim
  is a per-claim heuristic detection rate on a self-introduced benchmark, not
  a calibrated bound, so there is no coverage guarantee and no quantity a
  reader can hold the issuer to. The line this paper draws is therefore not
  detection versus silence about meaning; it is an unbound heuristic versus a
  publicly verifiable, distribution-free bound that replays offline.
\item \textbf{The semantic gap, named independently.} A 2026 survey of
  AI-content identity and provenance~\cite{survey2026} articulates exactly
  this gap --- that \emph{cryptographic correctness does not imply semantic
  correctness}, and that prevailing provenance attests output attribution,
  not ``whether the transformation preserved semantic fidelity.'' Our
  certificate is a concrete instrument for that independently-named gap.
\item \textbf{C2PA / Content Credentials}~\cite{c2pa} bind provenance to
  media; C2PA 2.4 adds text manifests, but the text binding is a
  \emph{byte-exact} hard hash that breaks under any edit or paraphrase, and
  the spec itself states provenance is not a truth or quality claim.
  \textbf{SynthID-Text} (Dathathri et al., 2024~\cite{dathathri2024}) and
  statistical detectors target generation, not
  transformation-preservation, and degrade under exactly the rewriting text
  invites --- a property re-confirmed in 2025 for SynthID
  specifically~\cite{synthidrobust2025} and across detector families by a
  single training-free paraphrase attack~\cite{advparaphrase2025};
  multi-hop rewriting~\cite{chainwash2026} drives watermark detection from
  $\sim$88\% to under 5\% while holding semantic similarity, the exact
  regime a meaning-preservation certificate is built for. We cite
  \cite{synthidrobust2025} in full rather than for its negative half only:
  the same work proposes SynGuard, a semantic-aware hybrid that recovers an
  average of 11.1\% F1 over SynthID-Text under those attacks. That is a
  genuine mitigation and the closest prior work in spirit. It remains
  detection: it raises the probability that a mark survives a rewrite, and
  makes no statement about what the rewrite preserved. The certificate here
  is orthogonal and composable with it, carrying an issuer-signed,
  replayable, distribution-free bound on a named meaning-loss proxy, which is
  a claim about the transformation rather than a probability about the
  artifact.
\item \textbf{SCITT (RFC 9943) and COSE Receipts (RFC 9942), June 2026.}
  The IETF's standards-track architecture for single-issuer signed statements
  with transparency receipts is now published~\cite{rfc9943}: an Issuer signs
  a Statement as a \texttt{COSE\_Sign1} Signed Statement, a Transparency
  Service registers it and returns a COSE Receipt (defined in RFC
  9942~\cite{rfc9942}), and the receipt travels in the statement's
  unprotected header as a Transparent Statement. The architecture states its
  own scope: ``Transparency does not prevent dishonest or compromised
  Issuers, but it holds them accountable.'' It certifies that a statement was
  made and registered; it makes no claim about the statement's content. That
  is complementary to this work, not competing: a
  \texttt{sum.meaning\_risk\_receipt.v1} is a natural Signed Statement
  payload, and registering one would add third-party witnessing of
  \emph{occurrence} to the issuer-signed bound on \emph{meaning}. We have
  not implemented the \texttt{COSE\_Sign1} re-envelope; the current wire
  format is JWS over JCS bytes (\S\ref{sec:substrate}), and the mapping is
  future work (\S\ref{sec:limits}).
\item \textbf{EU AI Act Article 50}~\cite{euaiact} obliges the
  \emph{generator} to mark; the associated \emph{Code of Practice on
  Transparency of AI-Generated Content} (finalized 10 June 2026) sets out
  the marking expectations. This work is an adjacent, complementary
  instrument (it attests what a \emph{transform} preserved), not a 50(2)
  marker substitute; we make no claim about which marking technologies the
  Code prescribes for text, and this certificate is complementary to
  whatever marker is used.
\item \textbf{Distribution-free guarantees.} Hoeffding
  (1963)~\cite{hoeffding1963}; the empirical-Bernstein bound (Maurer \&
  Pontil, 2009~\cite{maurerpontil2009}); conformal prediction (Vovk et al.,
  2005~\cite{vovk2005}; Angelopoulos \& Bates,
  2023~\cite{angelopoulosbates2023}) and conformal risk control
  (Angelopoulos et al., ICLR 2024~\cite{angelopoulos2024crc});
  group-conditional coverage (Gibbs, Cherian \& Cand\`es,
  2023~\cite{gibbs2023}). The nearest applications to \emph{transform
  quality} are conformal prediction over a machine-translation quality
  score~\cite{giovannotti2023} and conformal coverage on summary
  sentence-importance~\cite{confsumm2025}; the nearest
  \emph{output}-factuality line is Mohri \& Hashimoto~\cite{mohri2024} and
  conditional-factuality certificates from March 2026~\cite{condfact2026}. Two
  June 2026 results sit closest. \textbf{CARE}~\cite{care2026} applies
  conformal risk control to medical summarization, overlaying calibrated
  omission and hallucination flags with finite-sample, distribution-free
  guarantees; it is a post-hoc safety layer and produces no signed or
  verifiable artifact (the paper contains no notion of receipt, certificate,
  or provenance). \textbf{Kotte}~\cite{kotte2026} characterizes when
  conformal risk control can certify structured LLM outputs at all, proves an
  impossibility result, and analyzes a certification hierarchy across
  Hoeffding, empirical-Bernstein, and a betting-based e-CRC bound with strict
  gains in low-variance and large-sample regimes; that hierarchy is the one
  our \texttt{method} field names, and the betting bound is a candidate
  fourth rung we have not adopted. None of these couples its statistical
  bound to a signed, replayable receipt, and none bounds expected
  meaning-loss of a \emph{transform} under a named judge.
  Our contribution is the \emph{composition} --- a signed, replayable
  receipt over a named meaning-loss proxy --- not a new inequality.
\end{itemize}

\section{Limitations and future work}\label{sec:limits}

Validity rests on two assumptions separated in \S\ref{sec:certificate}:
independence within the calibration sample, which the three shipped
inequalities require, and a calibration-to-deployment match, which is assumed
rather than sampled and is not testable from the calibration set. Earlier
versions of this work, and the \texttt{disclosure} strings in the two
committed receipts, said ``exchangeability'' for both; that is the right
precondition for conformal \emph{prediction} but weaker than these
inequalities need, and the correction is stated rather than quietly applied. Model-judge replay is machine-pinned; de-pinning via
integer/fixed-point CPU inference (so a meaning-judge forward pass is
hardware-independent) is named future work --- same-hardware bitwise
determinism is now an engineering solved problem (batch-invariant kernels),
but cross-hardware reproducibility remains open, which is the niche this
de-pinning targets. The proxy is a proxy; stronger faithfulness judges are a
drop-in upgrade --- MiniCheck (Tang et al., 2024~\cite{minicheck2024}) is
the lightweight lineage anchor, and lighter open judges now report stronger
aggregate results, though on their own benchmarks: Paladin-mini (3.8B) reports
91.8\% average balanced accuracy against 78.2\% for Bespoke-MiniCheck-7B on
the grounding benchmark it introduces, while losing that benchmark's
time-and-date subset (82.0\% against 90.0\%)~\cite{paladinmini2025} --- and
any
judge-conditional bound inherits the judge's \emph{label-noise} ceiling, not
merely its accuracy. That last step is our inference, not a result of the work
we draw it from: Seo et al.~\cite{verifyingverifiers2025} report empirically
that approximately 16\% \emph{of} ambiguous or incorrectly labelled data
substantially influences model rankings on fact-verification benchmarks, which
is a statement about annotation quality rather than a flat benchmark error
rate, and we take the consequence for a judge-conditional bound from there. A tighter
betting/empirical-Bernstein confidence sequence is a no-wire-change
tightening. The corpora here demonstrate the mechanism; they are not a
benchmark suite. The binding open problem is not a sharper inequality but
\textbf{adoption}: one external party issuing and verifying a receipt it did
not author.

\section{Artifact availability}\label{sec:artifact}

Every receipt number in \S\ref{sec:billsum} and \S\ref{sec:translation} is
reproducible from committed bytes (the synthetic coverage sweep of
\S\ref{sec:coverage} is not: its generator is not committed, and it is
reported as a validation of the estimators rather than as an artifact),
and every verification claim in \S\ref{sec:verify} is executable by a third
party without contacting the authors.

\begin{itemize}
% Long, unhyphenatable typewriter paths in this list; \sloppy trades slightly
% loose interword space for keeping them out of the margin. itemize is a
% group, so this is scoped to this list only.
\sloppy
\item \textbf{Source.} \url{https://github.com/OtotaO/SUM}, Apache-2.0.
  Archived in Software Heritage for permanent citation as
  \texttt{swh:1:snp:1904bd383765f5cf553f08ea7ee46aca9925f25d} (origin
  \texttt{swh:1:ori:a7b5385a59a6fb561cbad53ce12da3149439401e}).
\item \textbf{Verifier.} \texttt{pip install "sum-engine[verify]"} installs
  the dependency-light \texttt{sum\_verify} SDK. The load-bearing promise is
  the exclusion, not a count: \textbf{no numpy, scipy, torch, GPU or network}
  on the verification path. It pulls \texttt{joserfc} plus the package's base
  \texttt{cryptography} and \texttt{sympy} (the latter is not imported by
  \texttt{sum\_verify}). To replay a committed golden with no arguments,
\begin{quote}\texttt{python -m sum\_verify -{}-demo}\end{quote}
  and to verify a receipt of your own, pass the files explicitly,
\begin{quote}\texttt{python -m sum\_verify <receipt> -{}-jwks <jwks.json> -{}-losses <losses.json>}\end{quote}
  which reproduces both stages of \S\ref{sec:verify}; the receipt alone is
  not enough, and \texttt{sum\_verify} says so rather than guessing.
  A second, independent implementation performs \textbf{Stage A} on the same
  bytes --- signature, schema gate, header invariants, disclosure --- namely
  the JS verifier
  \texttt{single\_file\_demo/\allowbreak meaning\_receipt\_\allowbreak verifier.js},
  run under Node in CI and in the browser demo. \textbf{Stage B replay is
  Python-only}: no JavaScript port of the bound kernel exists, so the
  cross-runtime claim in this paper covers authenticity and disclosure, not
  bound reproduction. (\texttt{standalone\_verifier/} verifies the older
  CanonicalBundle format, not meaning receipts.)
\item \textbf{Certificates.} Both binding-gate receipts reported in
  \S\ref{sec:empirical} are committed together with their public keys, their
  per-example loss vectors, and the generator that rebuilds them from the
  public corpus: \texttt{fixtures/\allowbreak meaning\_receipts\_\allowbreak billsum/}
  (BillSum, CC0) and
  \texttt{fixtures/\allowbreak meaning\_receipts\_\allowbreak translation/}
  (opus-100, EN to FR). Each directory contains the signed receipt,
  \texttt{jwks.json}, the losses file, and a
  \texttt{generate\_*\_fixture.py}.
\item \textbf{Scope of the artifact.} These artifacts replay the certificate
  and its bound from committed bytes. They do not re-run the judge model: the
  judge is named and pinned inside the receipt, and \S\ref{sec:boundary}
  states exactly what that dependence does and does not license a reader to
  conclude.
\end{itemize}

% References. The Markdown draft flags these as "to be formatted; primary
% sources verified". The two entries that once carried the draft's
% "[exact title to verify before submission]" note (survey2026,
% chainwash2026) were resolved against the live arXiv abstract pages on
% 2026-07-16 (PR #405); the note is now gone from this file and from the
% Markdown draft.
\begin{thebibliography}{99}

\bibitem{rfc8032}
S.~Josefsson and I.~Liusvaara.
\newblock {Edwards-Curve Digital Signature Algorithm (EdDSA)}.
\newblock RFC 8032, 2017.
\newblock (Ed25519; scheme due to D.~J. Bernstein et al.)

\bibitem{rfc8785}
A.~Rundgren, B.~Jordan, and S.~Erdtman.
\newblock {JSON Canonicalization Scheme (JCS)}.
\newblock RFC 8785, 2020.

\bibitem{rfc7515}
M.~Jones, J.~Bradley, and N.~Sakimura.
\newblock {JSON Web Signature (JWS)}.
\newblock RFC 7515, 2015.

\bibitem{rfc7517}
M.~Jones.
\newblock {JSON Web Key (JWK)}.
\newblock RFC 7517, 2015.

\bibitem{rfc9943}
H.~Birkholz, A.~Delignat-Lavaud, C.~Fournet, Y.~Deshpande, and S.~Lasker.
\newblock {An Architecture for Trustworthy and Transparent Digital Supply
Chains}.
\newblock RFC 9943 (Supply Chain Integrity, Transparency, and Trust),
Proposed Standard, June 2026.

\bibitem{rfc9942}
O.~Steele, H.~Birkholz, A.~Delignat-Lavaud, and C.~Fournet.
\newblock {CBOR Object Signing and Encryption (COSE) Receipts}.
\newblock RFC 9942, Proposed Standard, June 2026.

\bibitem{hoeffding1963}
W.~Hoeffding.
\newblock Probability inequalities for sums of bounded random variables.
\newblock \emph{Journal of the American Statistical Association},
  58(301):13--30, 1963.

\bibitem{maurerpontil2009}
A.~Maurer and M.~Pontil.
\newblock Empirical {B}ernstein bounds and sample variance penalization.
\newblock In \emph{COLT}, 2009.

\bibitem{vovk2005}
V.~Vovk, A.~Gammerman, and G.~Shafer.
\newblock \emph{Algorithmic Learning in a Random World}.
\newblock Springer, 2005.

\bibitem{angelopoulosbates2023}
A.~N. Angelopoulos and S.~Bates.
\newblock Conformal prediction: A gentle introduction.
\newblock \emph{Foundations and Trends in Machine Learning}, 2023.

\bibitem{angelopoulos2024crc}
A.~N. Angelopoulos et al.
\newblock Conformal risk control.
\newblock In \emph{ICLR}, 2024.

\bibitem{gibbs2023}
I.~Gibbs, J.~J. Cherian, and E.~J. Cand\`es.
\newblock Conformal prediction with conditional guarantees.
\newblock 2023.

\bibitem{dathathri2024}
S.~Dathathri et al.
\newblock Scalable watermarking for identifying large language model outputs
  ({SynthID-Text}).
\newblock \emph{Nature}, 2024.

\bibitem{aex2026}
Y.~Guan.
\newblock {AEX}: Non-intrusive multi-hop attestation and provenance for {LLM}
  {APIs}.
\newblock arXiv:2603.14283, 2026.

\bibitem{survey2026}
T.~Otsuka, K.~Toyoda, and A.~Leung.
\newblock {AI} identity: Standards, gaps, and research directions for {AI}
  agents.
\newblock arXiv:2604.23280, 2026.

\bibitem{chainwash2026}
M.~R. Ameen, A.~Islam, N.~Mahmud, and M.~E. Hamid.
\newblock Chainwash: Multi-step rewriting attacks on diffusion language model
  watermarks.
\newblock arXiv:2605.05503, 2026.

\bibitem{condfact2026}
K.~Ye, Q.~Pan, and S.~Li.
\newblock Conditional factuality controlled {LLMs} with generalization
  certificates via conformal sampling.
\newblock arXiv:2603.27403, 2026.

\bibitem{care2026}
S.~Bedi, B.~Lin, A.~Y. Zhou, C.~O. Stanwyck, J.~A. Jindal, S.~Koyejo,
D.~Stutz, and N.~H. Shah.
\newblock {CARE}: a conformal safety layer for medical summarization.
\newblock arXiv:2606.08969, 2026.

\bibitem{kotte2026}
V.~Kotte.
\newblock When can conformal risk control certify {LLM} outputs? {B}ounds,
impossibility, and adaptation for structured generation.
\newblock arXiv:2606.29054, 2026.

\bibitem{eigenai2026}
D.~Ribeiro~Alves, V.~Patankar, M.~Pereira, J.~Stephens, N.~Vaziri, and
S.~Kannan.
\newblock {EigenAI}: deterministic inference, verifiable results.
\newblock arXiv:2602.00182, 2026.

\bibitem{toolreceipts2026}
A.~Basu.
\newblock Tool receipts, not zero-knowledge proofs: practical hallucination
  detection for {AI} agents.
\newblock arXiv:2603.10060, 2026.

\bibitem{giovannotti2023}
P.~Giovannotti.
\newblock Evaluating machine translation quality with conformal predictive
  distributions.
\newblock arXiv:2306.01549 ({COPA} 2023).

\bibitem{confsumm2025}
B.~Kuwahara, C.-Y. Lin, X.~S. Huang, K.~K. Leung, J.~A. Yapeter,
I.~Stanevich, F.~Perez, and J.~C. Cresswell.
\newblock Document summarization with conformal importance guarantees.
\newblock {NeurIPS} 2025; arXiv:2509.20461.

\bibitem{mohri2024}
C.~Mohri and T.~Hashimoto.
\newblock Language models with conformal factuality guarantees.
\newblock arXiv:2402.10978, 2024.

\bibitem{synthidrobust2025}
X.~Han, Q.~Li, J.~Ni, and M.~Zulkernine.
\newblock Robustness assessment and enhancement of text watermarking for
  {Google's} {SynthID}.
\newblock arXiv:2508.20228, 2025.

\bibitem{advparaphrase2025}
Y.~Cheng, V.~S. Sadasivan, M.~Saberi, S.~Saha, and S.~Feizi.
\newblock Adversarial paraphrasing: a universal attack for humanizing
  {AI}-generated text.
\newblock {NeurIPS} 2025; arXiv:2506.07001.

\bibitem{minicheck2024}
L.~Tang et al.
\newblock {MiniCheck}: efficient fact-checking of {LLMs} on grounding
  documents.
\newblock In \emph{EMNLP}, 2024. arXiv:2404.10774.

\bibitem{paladinmini2025}
D.~Ivry and O.~Nahum.
\newblock Paladin-mini: a compact and efficient grounding model excelling in
real-world scenarios.
\newblock arXiv:2506.20384, 2025.

\bibitem{verifyingverifiers2025}
W.~Seo, S.~Han, J.~Jung, B.~Newman, S.~Lim, S.~Lee, X.~Lu, Y.~Choi, and
Y.~Yu.
\newblock Verifying the verifiers: unveiling pitfalls and potentials in fact
  verifiers.
\newblock {COLM} 2025; arXiv:2506.13342.

\bibitem{c2pa}
{C2PA}.
\newblock \emph{Coalition for Content Provenance and Authenticity, Technical
  Specification} (digitalSourceType taxonomy, v2.4; text manifests added in
  2.x).

\bibitem{euaiact}
European Union.
\newblock \emph{Artificial Intelligence Act}, Article 50 (applicable 2 Aug
  2026), and the \emph{Code of Practice on Transparency of AI-Generated
  Content} (finalized 10 June 2026).

\bibitem{ecma262}
{ECMA-262}.
\newblock \emph{ECMAScript Language Specification} (\S~Number::toString).

\end{thebibliography}

% ---------------------------------------------------------------------------
% Preserved meta-commentary from the Markdown draft (NOT paper content):
%
% Header note: "Draft (final-form v0) --- 2026-06-09. Targeted venue: cs.CR
% (primary), cs.LG (secondary). Every empirical number is cited from a
% committed, signed, independently re-derived receipt; every claim is scoped
% to the proof boundary (S8). Operator-review draft --- voice, emphasis, and
% final framing are the operator's to shape; drafting notes are at the foot."
%
% Drafting notes for the operator (next pass): (i) the Abstract section prose
% is 1,426 chars --- it fits the arXiv 1,920-char abstract field verbatim
% as-is (the "Contributions" list stays in the body; abstract + contributions
% together is 2,113 and would need cutting); (ii) decide whether S3's
% substrate (bench_digest reproducibility, the six compliance validators)
% earns its own subsection or stays compressed; (iii) add the coverage figure
% (the S7.3 table as a plot) and a system diagram; (iv) pick the single
% running example used end-to-end; (v) fetch full bibliographic details +
% DOIs and convert to the venue's bib style; (vi) reconcile first-person /
% voice. The sheaf-detector honest-negative results are intentionally absent
% --- they are the separate cs.LG Paper 2.
% ---------------------------------------------------------------------------

\end{document}
